\documentclass[a4paper,12pt]{report}

\PassOptionsToPackage{table,xcdraw}{xcolor}
\usepackage{tocloft}
\renewcommand{\cftchappresnum}{CHƯƠNG~}
\renewcommand{\cftchapaftersnum}{. }
\setlength{\cftchapnumwidth}{6em}

\usepackage[colorlinks=true, linkcolor=black, urlcolor=blue]{hyperref}
\usepackage[section]{placeins}
\usepackage{mathtools}
\usepackage{amssymb}
\usepackage{cleveref}
\usepackage{siunitx}
\usepackage{float}
\usepackage{graphicx}
\usepackage[justification=centering]{caption}
\usepackage{subcaption}
\usepackage{wallpaper}
\usepackage{nomencl}
\usepackage{fancyhdr}
\renewcommand{\chaptername}{CHƯƠNG}
\makeatletter
\renewcommand{\chaptermark}[1]{%
  \markboth{%
    \chaptername~\thechapter.\ #1%
  }{}%
}
\makeatother
\usepackage{url}
\usepackage[left=2.5cm,right=2.5cm,top=2cm,bottom=3.5cm]{geometry}

\newcommand{\UE}[1]{\renewcommand{\UE}{#1}}
\newcommand{\sujet}[1]{\renewcommand{\sujet}{#1}}
\newcommand{\titre}[1]{\renewcommand{\titre}{#1}}
\newcommand{\enseignant}[1]{\renewcommand{\enseignant}{#1}}
\newcommand{\eleves}[1]{\renewcommand{\eleves}{#1}}

\newcommand{\fairemarges}{
\makenomenclature
\pagestyle{fancy}
\fancyheadoffset{1cm}
\setlength{\headheight}{2cm}
\lhead{CS338.Q21 -- Báo cáo đồ án}
\rhead{\MakeUppercase{\leftmark}}
}

\newcommand{\fairepagedegarde}{
\begin{titlepage}
\IfFileExists{logos/background.jpg}{\ThisLRCornerWallPaper{1}{logos/background.jpg}}{}
    \centering

    {\fontsize{18}{18}\selectfont
    ĐẠI HỌC QUỐC GIA THÀNH PHỐ HỒ CHÍ MINH \\
    TRƯỜNG ĐẠI HỌC CÔNG NGHỆ THÔNG TIN
    \par}
    \vspace{1.5cm}
    {\scshape\Large \sujet \par}
    \vspace{1cm}
    \rule{\linewidth}{0.2 mm} \\[0.4 cm]
    {\huge\bfseries \titre \par}
    \rule{\linewidth}{0.2 mm} \\[1.5 cm]
    \vspace{1cm}

    \begin{minipage}{0.5\textwidth}
        \begin{flushleft} \large
        \emph{\textbf{Nhóm sinh viên:}}\\
        \eleves\\
        \end{flushleft}
    \end{minipage}
    ~
    \begin{minipage}{0.4\textwidth}
        \begin{flushright} \large
        \emph{\textbf{GVHD :}} \\
        \enseignant \\
        \end{flushright}
    \end{minipage}\\[2cm]

    \vfill
    {\large \today\par}
\end{titlepage}
}

\newcommand{\tabledematieres}{
\tableofcontents
\newpage
}

% Conditional figure include: drop in figures/<name> when available.
\newcommand{\figureorplaceholder}[3]{%
  \IfFileExists{#1}{%
    \includegraphics[width=#2]{#1}%
  }{%
    \fbox{\parbox{#2}{\centering \textit{[Placeholder]}\\\texttt{#1}\\\small #3}}%
  }%
}

\captionsetup{
    font=small,
    labelfont=bf
}

\numberwithin{figure}{chapter}
\numberwithin{table}{chapter}

\usepackage{lipsum}
\usepackage{xcolor}
\definecolor{bg}{rgb}{0.95,0.95,0.95}
\usepackage{listings}
\lstdefinestyle{bashstyle}{
    backgroundcolor=\color{bg},
    basicstyle=\ttfamily\small,
    breaklines=true,
    columns=fullflexible,
    keepspaces=true,
    language=bash,
    showstringspaces=false,
    frame=single,
    framesep=4pt,
    xleftmargin=4pt,
    xrightmargin=4pt
}
\usepackage{multirow}
\usepackage{adjustbox}
\usepackage{booktabs}
\usepackage{colortbl}
\usepackage{xurl}
\usepackage{fontspec}
\usepackage{polyglossia}
\renewcommand{\chaptername}{CHƯƠNG}
\setmainlanguage{vietnamese}
\setmainfont{texgyretermes}[
    Extension      = .otf,
    UprightFont    = *-regular,
    BoldFont       = *-bold,
    ItalicFont     = *-italic,
    BoldItalicFont = *-bolditalic
]
\usepackage{cite}
\usepackage{pgfplots}
\usepackage{pgfplotstable}
\pgfplotsset{compat=1.18}

\renewcommand{\bibname}{Tài liệu tham khảo}
\setcounter{tocdepth}{3}
\setcounter{secnumdepth}{3}

\addto\captionsvietnamese{%
  \renewcommand{\contentsname}{MỤC LỤC}
  \renewcommand{\listfigurename}{DANH SÁCH HÌNH ẢNH}
  \renewcommand{\listtablename}{DANH SÁCH BẢNG}
}

\begin{document}

\titre{Đề tài: \\
Tên tiếng Việt: Đếm đối tượng Zero-shot với mẫu chất lượng \\
Tên tiếng Anh: Zero-shot Object Counting with Good Exemplars}
\sujet{NHẬN DẠNG -- CS338.Q21 \\ BÁO CÁO ĐỒ ÁN}
\UE{ }
\enseignant{TS. Dương Việt Hằng}
\eleves{Nguyễn Công Phát -- 23521143 (Nhóm trưởng) \\
Mai Thái Bình -- 23520158 \\
Vũ Việt Cường -- 23520213}

\fairemarges
\fairepagedegarde

\tabledematieres

\listoffigures
\vspace{1cm}

\listoftables
\newpage

%------------ LỜI NÓI ĐẦU ----------------
\chapter*{LỜI NÓI ĐẦU}
\addcontentsline{toc}{chapter}{LỜI NÓI ĐẦU}

Lời đầu tiên, nhóm chúng em xin bày tỏ lòng biết ơn sâu sắc đến cô \textbf{TS. Dương Việt Hằng}, người đã tận tình hướng dẫn, định hướng và truyền đạt những kiến thức quý báu trong môn học \textit{Nhận dạng (CS338.Q21)}. Sự góp ý chuyên môn và những buổi trao đổi của cô đã giúp nhóm hiểu rõ hơn về bài toán đếm đối tượng zero-shot, hệ thống VA-Count cũng như cách đánh giá khoa học một mô hình thị giác máy tính trên bộ dữ liệu FSC147.\\[1em]
Đồ án này hiện thực lại hệ thống VA-Count~\cite{vacount} từ bài báo gốc và mở rộng theo hai hướng độc lập trong cùng dòng tài liệu: \textit{Rich Prompts}~\cite{richprompts} cho bước sinh exemplar và \textit{YOLO-World}~\cite{yoloworld} thay cho GroundingDINO trong khâu phát hiện. Trên cơ sở đó, nhóm CS338.Q21 chúng em đã tổ chức kho mã nguồn (\texttt{docs/}, \texttt{configs/}, \texttt{scripts/}), chuẩn hóa toàn bộ bình luận sang tiếng Anh, đưa khóa Gemini API ra biến môi trường, viết tài liệu và soạn báo cáo này dưới góc nhìn của bài toán nhận dạng. Toàn bộ kết quả thực nghiệm được lấy từ log \texttt{wandb} của chính nhóm, lưu trong thư mục
\mbox{\texttt{experiments/exp\{2,3,4,5\}/wandb/}}
và được dẫn nguồn rõ ràng tại Chương~\ref{ch:experiments}.\\[1em]
Chúng em cũng xin chân thành cảm ơn các bạn lớp CS338.Q21 đã cùng thảo luận, kiểm thử và đóng góp ý kiến trong suốt quá trình thực hiện. Nhóm rất mong nhận được thêm góp ý từ cô để hoàn thiện kho mã nguồn cũng như báo cáo này.\\[1em]
\begin{flushright}
\textit{Thành phố Hồ Chí Minh, \today}
\end{flushright}

%------------ THÔNG TIN THÀNH VIÊN ----------------
\chapter*{THÔNG TIN THÀNH VIÊN}
\addcontentsline{toc}{chapter}{THÔNG TIN THÀNH VIÊN}

Bảng~\ref{tab:phan_cong} trình bày phân công nhiệm vụ trong nhóm CS338.Q21 đối với đồ án \textit{“Zero-shot Object Counting with Good Exemplars”}.

\begin{table}[H]
\centering
\renewcommand{\arraystretch}{1.25}
\begin{tabular}{|l|c|c|l|p{6cm}|}
\hline
\textbf{Họ và tên} & \textbf{MSSV} & \textbf{Đóng góp} & \textbf{Vai trò} & \textbf{Phần việc đảm nhận} \\
\hline
Nguyễn Công Phát & 23521143 & 35\% & Nhóm trưởng &
Tổ chức kho mã nguồn (\texttt{docs/}, \texttt{configs/}, \texttt{scripts/}); chuẩn hóa toàn bộ bình luận sang tiếng Anh; đưa \texttt{GEMINI\_API\_KEY} ra \texttt{.env}; tích hợp pipeline huấn luyện và đánh giá (\texttt{FSC\_train.py}, \texttt{FSC\_test.py}); viết Chương~1 và phần phương pháp. \\
\hline
Mai Thái Bình & 23520158 & 35\% & Thành viên &
Hiện thực Module Rich Prompt (Gemini 2.0 Flash) cùng pipeline \texttt{grounding\_pos.py} / \texttt{grounding\_neg.py}; chạy các thử nghiệm đánh giá MAE/RMSE; thiết kế bảng kết quả và biểu đồ tốc độ--độ chính xác; viết Chương~2 (Rich Prompt) và phần phân tích lỗi sai. \\
\hline
Vũ Việt Cường & 23520213 & 30\% & Thành viên &
Hiện thực mở rộng với YOLO-World (\texttt{yolo\_pos\_withPrompt.py}, \texttt{yolo\_pos\_withoutPrompt.py}, \texttt{yolo\_neg.py}); kiểm thử và viết tài liệu cho ứng dụng demo Streamlit (\texttt{demo\_app\_advanced.py}) cùng các module hỗ trợ; thiết kế slide báo cáo; viết Chương~3 (YOLO-World), Chương~4 (Demo) và Chương~5. \\
\hline
\multicolumn{5}{|l|}{\textit{Tổng cộng: 100\%.}} \\
\hline
\end{tabular}
\caption{Phân chia công việc và mức đóng góp của các thành viên CS338.Q21}
\label{tab:phan_cong}
\end{table}

Mức độ đóng góp đã được các thành viên thảo luận, thống nhất và đánh giá khách quan dựa trên khối lượng công việc thực tế đã hoàn thành.

%==================================================================
%                          CHƯƠNG 1
%==================================================================
\chapter{TỔNG QUAN ĐỀ TÀI}
\label{ch:overview}

\section{Mô tả bài toán}
\textbf{Đếm đối tượng (Object Counting)} là bài toán cơ bản trong thị giác máy tính, yêu cầu hệ thống ước lượng số lượng các thực thể thuộc một lớp ngữ nghĩa cho trước trong một ảnh tự nhiên. Khác với các tiếp cận \textit{class-specific} (chỉ đếm một lớp được huấn luyện sẵn) hay \textit{few-shot} (yêu cầu người dùng cung cấp một vài ảnh mẫu), nhánh \textbf{Zero-shot Counting} đặt mục tiêu loại bỏ hoàn toàn dữ liệu mẫu hình ảnh được cung cấp thủ công: hệ thống chỉ nhận một ảnh đầu vào và một mô tả văn bản (\textit{text prompt}) về lớp cần đếm, sau đó tự động xác định mẫu (exemplars) và sinh \textit{density map} để tính tổng số lượng đối tượng~\cite{vacount, richprompts}.

\section{Lý do chọn đề tài}
Sự bùng nổ của dữ liệu thị giác trong kỷ nguyên số đặt ra nhu cầu cấp thiết về các hệ thống phân tích hình ảnh tự động, trong đó đếm đối tượng đóng vai trò then chốt cho các ứng dụng nông nghiệp thông minh, giám sát giao thông, kiểm kê kho hàng, và phân tích hình ảnh y tế. Tuy nhiên, các tiếp cận hiện nay vẫn tồn tại nhiều hạn chế:

\begin{itemize}
    \item \textbf{Phụ thuộc vào nhãn thủ công.} Các phương pháp giám sát yêu cầu lượng lớn nhãn dạng \textit{bounding box} hoặc \textit{dot annotations}, vốn tốn kém và không khả thi trên dữ liệu quy mô lớn.
    \item \textbf{Hạn chế miền đối tượng.} Các mô hình giám sát thường suy giảm đáng kể trên các lớp chưa từng xuất hiện trong huấn luyện (\textit{unseen classes}).
    \item \textbf{Tính thủ công của Few-shot.} Người dùng vẫn phải tự cung cấp ảnh mẫu, cản trở khả năng tự động hóa.
\end{itemize}

Đồ án này hướng tới ba mục tiêu chính:

\begin{enumerate}
    \item \textbf{Tự động hóa quy trình lấy mẫu:} loại bỏ hoàn toàn sự can thiệp của con người trong việc chọn exemplars.
    \item \textbf{Học cơ chế liên kết đa phương thức:} xây dựng tương tác giữa ngữ nghĩa văn bản (text prompts) và đặc trưng hình ảnh.
    \item \textbf{Tối ưu hóa khả năng tổng quát hóa:} duy trì độ chính xác trên các lớp đối tượng mới mà mô hình chưa từng nhìn thấy trong quá trình huấn luyện.
\end{enumerate}

\section{Phát biểu bài toán}
\label{sec:problem-statement}

\paragraph{Đầu vào.}
\begin{itemize}
    \item Một ảnh màu RGB $I \in \mathbb{R}^{H \times W \times 3}$ chứa các đối tượng trong bối cảnh tự nhiên.
    \item Một văn bản (text prompt) $T$ là tên lớp tiếng Anh mô tả đối tượng cần đếm (ví dụ: \texttt{"oranges"}, \texttt{"keyboard keys"}).
    \item Tập dữ liệu chuẩn FSC147~\cite{ranjan2021fsc147} dùng để huấn luyện và đánh giá.
\end{itemize}

\paragraph{Đầu ra.}
\begin{itemize}
    \item Một bản đồ mật độ $D \in \mathbb{R}^{H \times W}$ biểu diễn phân bố không gian của các đối tượng.
    \item Tổng số lượng đối tượng $\hat{y} = \int_{H,W} D \, dx\, dy$ tương ứng với mô tả $T$.
\end{itemize}

\paragraph{Ràng buộc.}
\begin{itemize}
    \item \textbf{Dữ liệu ảnh:} ảnh RGB có thể chứa một hoặc nhiều đối tượng với mật độ rất biến thiên; mỗi lần thực thi chỉ đếm \emph{một lớp} đối tượng.
    \item \textbf{Câu lệnh:} bằng tiếng Anh; mô tả vật thể hữu hình quan sát được; không sử dụng khái niệm trừu tượng (\texttt{traffic}, \texttt{meeting}) hay danh từ không đếm được (\texttt{water}, \texttt{smoke}).
\end{itemize}

\section{Lý do lựa chọn VA-Count làm baseline}
\label{sec:why-vacount}

Trong số các phương pháp zero-shot counting hiện tại, nhóm lựa chọn VA-Count~\cite{vacount} làm nền tảng vì ba lý do chính:

\begin{enumerate}
    \item \textbf{Cơ chế tự sinh exemplar hoàn toàn tự động.} Khác với CounTR~\cite{countr} hay FamNet~\cite{famnet} vốn yêu cầu người dùng cung cấp bounding box mẫu, VA-Count sử dụng GroundingDINO~\cite{groundingdino} để tự động sinh exemplar từ text prompt. Đây là tiền đề then chốt để hệ thống hoạt động theo chế độ zero-shot thực sự, phù hợp với mục tiêu loại bỏ hoàn toàn sự can thiệp thủ công.
    \item \textbf{Kiến trúc hai luồng đối lập (contrastive dual-branch).} Các phương pháp counting trước đó như BMNet~\cite{bmnet} chỉ sử dụng mẫu dương để sinh density map, dẫn đến dễ nhầm lẫn giữa đối tượng cần đếm và các vật thể nền có đặc trưng thị giác tương tự. VA-Count giải quyết vấn đề này bằng cách bổ sung một nhánh âm (negative branch) với cơ chế học tương phản, giúp mô hình phân biệt rõ ràng giữa foreground và background. Cơ chế contrastive learning này được chứng minh hiệu quả trong nhiều bài toán thị giác~\cite{simclr} và đặc biệt phù hợp với bài toán counting nơi nhiễu nền là thách thức lớn.
    \item \textbf{Thiết kế module hóa, dễ mở rộng.} Kiến trúc VA-Count tách biệt rõ ràng giữa pha lấy mẫu (EEM) và pha sinh density map (NSM). Điều này cho phép thay thế từng thành phần một cách độc lập: có thể thay GroundingDINO bằng YOLO-World trong EEM mà không cần thay đổi NSM, hoặc bổ sung module sinh prompt mà không ảnh hưởng đến Counter. Tính module hóa này là lợi thế lớn so với các kiến trúc end-to-end như ZSC~\cite{zeroshot_counting_xu}.
\end{enumerate}

Hình~\ref{fig:vis_posneq} minh họa cách hệ thống VA-Count hoạt động trên thực tế: mẫu dương (bounding box xanh lá) xác định đúng đối tượng cần đếm (strawberries), trong khi mẫu âm (bounding box đỏ) đánh dấu các vùng nền. Sự phối hợp giữa hai luồng giúp density map tập trung chính xác vào đối tượng mục tiêu.

\begin{figure}[H]
    \centering
    \figureorplaceholder{figures/vis_strawberry_posneq.jpg}{0.95\textwidth}{Minh họa VA-Count trên ảnh \texttt{346.jpg}: bounding box xanh lá = positive exemplars, bounding box đỏ = negative exemplars. Overlay density map cho thấy mô hình tập trung đúng vào các quả dâu tây.}
    \caption{Minh họa hoạt động của VA-Count: positive exemplars (xanh lá) và negative exemplars (đỏ) trên ảnh \texttt{346.jpg} (lớp \emph{strawberries}). Phần overlay bên phải cho thấy density map tập trung đúng vào đối tượng cần đếm.}
    \label{fig:vis_posneq}
\end{figure}

\section{Đóng góp của đồ án}
Đồ án này hiện thực lại bài báo VA-Count~\cite{vacount} từ đầu và mở rộng theo hai hướng độc lập trong cùng dòng tài liệu. Các đóng góp chính:

\begin{enumerate}
    \item \textbf{Hiện thực hệ thống VA-Count cơ sở.} Module \emph{Exemplar Enhancement Module} (EEM) sử dụng GroundingDINO~\cite{groundingdino} để sinh các hộp giới hạn ứng viên cho lớp đối tượng cần đếm; module \emph{Noise Suppression Module} (NSM) học mất mát đối lập giữa exemplar dương và exemplar âm theo công thức trong~\cite{vacount}.
    \item \textbf{Tích hợp Rich Prompts cho khâu lấy mẫu.} Áp dụng đề xuất của~\cite{richprompts}: dùng Gemini sinh mô tả phong phú từ tên lớp đầu vào, sau đó dùng CLIP~\cite{clip} re-rank các đề xuất để chọn ra ảnh mẫu chất lượng cao hơn.
    \item \textbf{Tích hợp YOLO-World thay cho GroundingDINO.} Sử dụng mô hình phát hiện vốn từ vựng mở YOLO-World~\cite{yoloworld} làm bộ phát hiện thay thế trong giai đoạn EEM, cho thời gian sinh exemplar nhanh hơn nhiều lần (xem Chương~\ref{ch:experiments}).
    \item \textbf{Đánh giá đầy đủ trên FSC-147.} Huấn luyện lại từ checkpoint backbone MAE~\cite{mae} và đánh giá MAE / RMSE / thời gian suy luận trên cả bốn cấu hình (baseline, +Rich Prompt, +YOLO-World, +YOLO-World+Rich Prompt); toàn bộ log lưu trong thư mục \mbox{\texttt{experiments/exp\{2,3,4,5\}/wandb/}}.
    \item \textbf{Báo cáo và demo.} Viết báo cáo song ngữ 26 trang dưới góc nhìn của bài toán nhận dạng và xây dựng ứng dụng demo Streamlit cho phép thử nghiệm hệ thống trực tiếp.
\end{enumerate}

Hai chi tiết cài đặt sau được nêu lại để bạn đọc tiện đối chiếu với hai bài báo gốc khi đọc Chương~\ref{ch:method}: (i) phiên bản Gemini sử dụng tại thời điểm cài đặt là \texttt{gemini-2.0-flash-exp}~\cite{gemini} (Rich Prompts trong~\cite{richprompts} đề cập đến Gemini~2.5~Flash); (ii) hàm mất mát của NSM được hiện thực theo dạng \emph{margin-ReLU contrastive} $\mathcal{L}_{C}=\mathrm{ReLU}(\ell_p-\ell_n+m)$ với $m{=}0{.}5$ (xem \texttt{FSC\_train.py}), tương đương về mặt tối ưu với dạng softmax trong~\cite{vacount} nhưng ổn định số học hơn khi huấn luyện trên FSC-147.

%==================================================================
%                          CHƯƠNG 2
%==================================================================
\chapter{PHƯƠNG PHÁP THỰC HIỆN}
\label{ch:method}

Chương này trình bày kiến trúc hệ thống mà nhóm hiện thực, gồm ba phần: (i) baseline \textit{VA-Count}~\cite{vacount}, (ii) mở rộng theo \textit{Rich Prompts}~\cite{richprompts} dùng LLM cho khâu sinh exemplar, và (iii) thay GroundingDINO bằng \textit{YOLO-World}~\cite{yoloworld} ở module phát hiện. Mỗi mục đều ghi rõ chi tiết cài đặt thực tế (phiên bản mô hình, hàm mất mát, siêu tham số) để bảo đảm tính tái lập.

\section{Kiến trúc tổng quan của VA-Count}
\label{sec:vacount-overview}

Sau khi nhận đầu vào (ảnh + class name), hệ thống chia thành hai luồng xử lý song song:

\begin{itemize}
    \item \textbf{Luồng dương (Positive):} tìm các mẫu dương đại diện cho đối tượng cần đếm. Kết quả là bản đồ mật độ dương $D_p$ tập trung điểm sáng tại đối tượng cần đếm.
    \item \textbf{Luồng âm (Negative):} tìm các mẫu âm đại diện cho vật thể nhiễu. Kết quả là bản đồ mật độ âm $D_n$ tập trung điểm sáng tại các vật thể không phải mục tiêu.
\end{itemize}

Hệ thống dùng hai khối phối hợp:

\begin{enumerate}
    \item \textbf{Exemplar Enhancement Module (EEM):} tự động phát hiện và tinh chỉnh exemplars.
    \item \textbf{Noise Suppression Module \& Counter (NSM):} sinh density map và khử nhiễu.
\end{enumerate}

\begin{figure}[H]
    \centering
    \figureorplaceholder{figures/system_overview.png}{0.95\textwidth}{Tổng quan kiến trúc VA-Count: hai luồng song song positive/negative qua EEM và NSM.}
    \caption{Tổng quan kiến trúc VA-Count gồm hai luồng positive/negative qua EEM và NSM. Sơ đồ chuyển vẽ theo kiến trúc đề xuất trong VA-Count~\cite{vacount}.}
    \label{fig:system_overview}
\end{figure}

\subsection{Exemplar Enhancement Module (EEM)}
\label{subsec:eem}
EEM có ba bước:

\paragraph{Bước 1: Khởi tạo mẫu bằng GroundingDINO~\cite{groundingdino}.}
\begin{itemize}
    \item \textit{Luồng dương:} truyền ảnh + class name (ví dụ \texttt{single dog .}) để GroundingDINO phát hiện sơ bộ các bounding box ứng viên dương.
    \item \textit{Luồng âm:} truyền ảnh kèm prompt chung (\texttt{object .} hoặc top-5 lớp âm trong dataset, xem \texttt{grounding\_neg.py}). GroundingDINO trả về tất cả vật thể có thể có trong ảnh, bao gồm cả mục tiêu và nhiễu.
\end{itemize}

\paragraph{Bước 2: Tinh chỉnh và lọc mẫu chất lượng.}
\begin{itemize}
    \item \textbf{Deduplication (chỉ luồng âm).} So sánh box âm với box dương: nếu IoU cao thì loại, đảm bảo mẫu âm không trùng mẫu dương.
    \item \textbf{Single-Object Filter.} Cả hai luồng đi qua bộ phân loại nhị phân $\phi_{\text{single}}$ dựa trên \textit{CLIP ViT-B/32} (xem \texttt{biclassify.py}). Patch chỉ được giữ khi $P(\phi_{\text{single}}(x)=1) > 0.8$.
\end{itemize}

\paragraph{Bước 3: Chọn top-$k$ mẫu.} Mặc định: top-3 cho mỗi luồng theo confidence score của GroundingDINO. Phiên bản Rich Prompt mở rộng lên top-5 cho luồng dương (Mục~\ref{sec:rich-prompts}).

\subsection{Noise Suppression Module \& Counter (NSM)}
\label{subsec:nsm}
NSM phân biệt rõ giữa đối tượng cần đếm và vật thể nhiễu, dựa trên các mẫu đã chọn từ EEM. Có hai pha chính:

\paragraph{Pha 1: Tương tác đặc trưng và sinh density map.}
Hệ thống dùng một bộ đếm gồm \textit{Image Encoder}, \textit{Interaction Module} và \textit{Decoder}:
\begin{itemize}
    \item Đặc trưng ảnh đầu vào đóng vai trò \textit{Query}.
    \item Phép chiếu tuyến tính của đặc trưng mẫu (cả $E_p$ và $E_n$) đóng vai trò \textit{Key/Value}.
\end{itemize}
Sau cơ chế attention và Decoder, hệ thống sinh hai bản đồ mật độ:
\begin{equation}
D_p = \Gamma_{\text{decode}}(F^p_{\text{fuse}}), \qquad
D_n = \Gamma_{\text{decode}}(F^n_{\text{fuse}}).
\end{equation}

\paragraph{Pha 2: Hàm mất mát.}
Bài báo VA-Count gốc~\cite{vacount} (Công thức (2)) trình bày hàm contrastive dạng softmax. Trong cài đặt của nhóm, \texttt{FSC\_train.py} dùng một hàm tương phản dạng \emph{margin-ReLU} cộng thẳng với mất mát mật độ. Dạng này tương đương về mặt tối ưu nhưng ổn định số học hơn khi huấn luyện trên FSC-147:

\begin{equation}
\mathcal{L}_{\text{pos}} = \frac{1}{HW}\sum_{i,j} M_{ij}\,(D^{(i,j)}_p - D^{(i,j)}_g)^2,
\qquad
\mathcal{L}_{\text{neg}} = \frac{1}{HW}\sum_{i,j} M_{ij}\,(D^{(i,j)}_n - D^{(i,j)}_g)^2,
\end{equation}
\begin{equation}
\mathcal{L}_{C} = \mathrm{ReLU}\!\left(\mathcal{L}_{\text{pos}} - \mathcal{L}_{\text{neg}} + m\right),\quad m = 0.5,
\qquad
\mathcal{L}_{\text{total}} = \mathcal{L}_{C} + \mathcal{L}_{\text{pos}}.
\end{equation}
Trong đó $M$ là mặt nạ Bernoulli ngẫu nhiên với $p=0{.}8$ áp lên mỗi pixel. Công thức softmax gốc trong~\cite{vacount} được giữ làm tham chiếu lý thuyết; cài đặt được dùng cho mọi kết quả trong báo cáo này là dạng margin-ReLU ở trên.

\section{Cải tiến với Rich Prompts}
\label{sec:rich-prompts}

\subsection{Động lực cải tiến}

Trong baseline VA-Count, bộ phát hiện GroundingDINO~\cite{groundingdino} nhận đầu vào là tên lớp đơn giản (ví dụ \texttt{"strawberries ."}). Cách tiếp cận này bộc lộ hai hạn chế quan trọng:

\begin{itemize}
    \item \textbf{Nhập nhằng ngữ nghĩa (semantic ambiguity).} Một class name ngắn có thể khớp với nhiều loại đối tượng khác nhau: ví dụ \texttt{"key"} có thể là phím bàn phím, chìa khóa, hoặc nốt nhạc. Nghiên cứu về prompt engineering trong mô hình vision-language~\cite{coop} chỉ ra rằng các mô tả chi tiết (rich descriptions) giúp giảm đáng kể sự nhập nhằng này trong không gian embedding.
    \item \textbf{Thiếu thông tin phân biệt foreground/background.} Khi ảnh có nền phức tạp (nhiều vật thể cùng màu, cùng hình dạng), một class name đơn không cung cấp đủ đặc trưng để GroundingDINO phân biệt đối tượng cần đếm với các vật thể gây nhiễu. Đây là vấn đề đã được nhận diện trong các nghiên cứu về open-vocabulary detection~\cite{ovdet_survey}.
\end{itemize}

Bài báo Rich Prompts~\cite{richprompts} giải quyết cả hai vấn đề bằng cách sử dụng một LLM đa phương thức (Multimodal LLM) để sinh prompt mô tả chi tiết đặc trưng thị giác cốt lõi của đối tượng. Nhóm hiện thực ý tưởng này với \emph{Gemini~2.0~Flash} (\texttt{gemini-2.0-flash-exp})~\cite{gemini}, là phiên bản có sẵn ở thời điểm cài đặt. Bài báo gốc dùng Gemini~2.5~Flash, và đây là một lựa chọn cài đặt khác biệt nhưng tương đương về vai trò sinh enhanced prompt.

\subsection{Module Enhanced Prompt}
\label{subsec:enhanced-prompt}

Module nhận ba loại đầu vào:
\begin{itemize}
    \item \textbf{Class name:} tên lớp đối tượng.
    \item \textbf{Image:} ảnh gốc.
    \item \textbf{System prompt:} gồm hai mẫu, đó là \textit{Positive Prompt Template} (mô tả đặc trưng thị giác cốt lõi) và \textit{Negative Prompt Template} (mô tả nền/nhiễu).
\end{itemize}

Chi tiết các prompt mẫu được thiết kế để định hướng mô hình tập trung vào
đặc trưng thị giác cốt lõi:

\begin{lstlisting}[style=bashstyle,caption={Mẫu Positive Prompt (Trích xuất đặc trưng)},label=lst:pos_prompt]
Look at the image and provide the visual definition
of a single '{singular_name}'.

Task: Describe the intrinsic physical appearance of
just ONE instance, as if it were isolated.

STRICT RULES:
1. Subject: Describe ONE item only.
   Ignore the background or other identical items.
2. Format: Start with 'single {singular_name}'.
   Use dot-separated phrases.
3. Content: Focus on Shape, Color, Material.

Example for 'keyboard key':
BAD:  keyboard key . rows of buttons . full layout .
GOOD: single keyboard key . square shape .
      black plastic material . white printed letter .

Your output for '{singular_name}':
\end{lstlisting}

\begin{lstlisting}[style=bashstyle,caption={Mẫu Negative Prompt (Loại bỏ nhiễu)},label=lst:neg_prompt]
Analyze the image and identify visual elements that
are NOT the object '{class_name}'.

Describe the background, environment, and potential
distractors.

Focus strictly on:
1. Background materials.
2. Environmental conditions.
3. Other objects that are NOT '{class_name}'.

Output format rules:
- Do NOT describe the '{class_name}' itself.
- Separate phrases with ' . '.
- Do NOT write full sentences.
- End with a dot.

Example input:  Image of a cat on a sofa.
Example output: beige sofa texture . blurry living
               room . shadows . fabric patterns .

Your output for this image:
\end{lstlisting}

\begin{figure}[H]
    \centering
    \figureorplaceholder{figures/enhanced_prompt_flow.png}{0.85\textwidth}{Sơ đồ luồng hoạt động của Module Enhanced Prompt: Image+ClassName+SystemPrompt → Gemini → enhanced positive/negative prompts.}
    \caption{Sơ đồ luồng hoạt động của Module Enhanced Prompt. Sơ đồ chuyển vẽ theo thiết kế trong Rich Prompts~\cite{richprompts}.}
    \label{fig:enhanced_prompt}
\end{figure}

\begin{figure}[H]
    \centering
    \begin{subfigure}{0.42\textwidth}
        \centering
        \figureorplaceholder{figures/positive_keyboard.png}{\linewidth}{Kết quả top-5 positive boxes cho lớp \texttt{keyboard keys} sau Re-ranking bằng CLIP ViT-L/14.}
        \caption{Class name: keyboard keys}
    \end{subfigure}\hfill
    \begin{subfigure}{0.42\textwidth}
        \centering
        \figureorplaceholder{figures/positive_grapes.png}{\linewidth}{Kết quả top-5 positive boxes cho lớp \texttt{grapes}.}
        \caption{Class name: grapes}
    \end{subfigure}
    \caption{Trích xuất Positive Boxes từ Module Enhanced Prompt. Hai cặp ảnh (trái/phải): bên trái là ảnh gốc + class name; bên phải là enhanced prompt do Gemini sinh ra cùng các Positive Boxes thu được. Minh họa theo cách trình bày trong Rich Prompts~\cite{richprompts}.}
    \label{fig:positive_boxes}
\end{figure}

\subsection{Kiến trúc cải tiến}

Hai luồng được sửa độc lập:

\textbf{Luồng dương:}
\begin{itemize}
    \item \textit{Input Query.} Dùng \emph{positive enhanced prompt} thay vì class name đơn thuần.
    \item \textit{Re-ranking.} Tích hợp \emph{CLIP ViT-L/14}~\cite{clip} đánh giá lại tương đồng giữa class name gốc và các patches do GroundingDINO đề xuất. Đây là sai khác đáng chú ý so với mô tả ban đầu (ViT-B/32); cụ thể, mã nguồn \texttt{grounding\_pos.py} dòng 49 dùng \texttt{clip.load("ViT-L/14")}.
    \item \textit{Chiến lược.} Thay vì top-3 như VA-Count gốc, lấy \textbf{top-5} patches có CLIP score cao nhất.
\end{itemize}

\textbf{Luồng âm:}
\begin{itemize}
    \item \textit{Input Query.} Thử nghiệm \emph{negative enhanced prompt} thay query mặc định \texttt{"object ."}.
    \item \textit{Deduplication.} Vẫn giữ class name (không phải enhanced prompt) làm tham chiếu để tránh triệt tiêu nhầm các vùng background hữu ích.
\end{itemize}

\begin{figure}[H]
    \centering
    \figureorplaceholder{figures/system_richprompt.png}{0.95\textwidth}{Kiến trúc VA-Count + Rich Prompt: Gemini sinh prompt → GroundingDINO → CLIP ViT-L/14 re-rank → NSM.}
    \caption{Kiến trúc hệ thống sau khi tích hợp Rich Prompt và CLIP Re-ranking. Sơ đồ chuyển vẽ theo thiết kế đề xuất trong Rich Prompts~\cite{richprompts}.}
    \label{fig:system_richprompt}
\end{figure}

\paragraph{Lý giải lựa chọn kỹ thuật.}
\begin{enumerate}
    \item \textbf{CLIP Re-ranking là cần thiết.} Khi dùng enhanced prompt, GroundingDINO trở nên ``nhạy'' hơn và trích xuất cả các vật có đặc tính tương tự nhưng không phải mục tiêu; CLIP đóng vai trò bộ lọc ngữ nghĩa cấp 2.
    \item \textbf{Deduplication theo class name.} Tránh tình trạng triệt tiêu hoàn toàn negative boxes nếu dùng enhanced prompt (vốn sinh quá nhiều positive matches).
    \item \textbf{Top-3 → Top-5.} Cung cấp thêm ngữ cảnh cho Image Encoder để cân bằng lại noise do số lượng candidates tăng.
\end{enumerate}

\section{Cải tiến với YOLO-World}
\label{sec:yolo-world}

\subsection{Động lực thay thế GroundingDINO}

Trong baseline VA-Count, GroundingDINO~\cite{groundingdino} là nút thắt cổ chai về tốc độ. Với kiến trúc Transformer kết hợp fusion cross-modal, mô hình cần tính toán Vision-Language Fusion cho mỗi khung hình trong thời gian thực, dẫn đến độ trễ cao và tiêu tốn tài nguyên đáng kể. Cụ thể, trong pipeline VA-Count, bước trích xuất exemplar chiếm tới $\sim$80\% tổng thời gian xử lý mỗi ảnh.

Vấn đề này đặc biệt nghiêm trọng khi triển khai thực tế:
\begin{itemize}
    \item \textbf{Thời gian chuẩn bị dữ liệu.} Trên toàn bộ 6.146 ảnh FSC-147, GroundingDINO tốn $\approx$13 giờ cho cả hai luồng (positive + negative). Con số này tăng lên 25 giờ khi bật Rich Prompt. Đây là rào cản lớn cho việc thử nghiệm nhanh các cấu hình khác nhau.
    \item \textbf{Ứng dụng tương tác.} Khi chạy demo, người dùng cần chờ $\sim$1,5 giây/ảnh chỉ cho bước phát hiện, chưa kể bước sinh density map. Yêu cầu phản hồi nhanh trong ứng dụng thực tế đòi hỏi thời gian xử lý dưới 1 giây tổng.
\end{itemize}

\subsection{Lựa chọn YOLO-World}

Nhóm thay bộ phát hiện bằng \emph{YOLO-World}~\cite{yoloworld}, một mô hình mở từ vựng nhẹ. Theo bài báo gốc, YOLO-World nhanh hơn GroundingDINO khoảng 20$\times$ nhờ chiến lược \textit{Prompt-then-Detect} và kiến trúc RepPAN (Re-parameterizable Path Aggregation Network). Cốt lõi của chiến lược này nằm ở việc mã hóa trước (offline) các gợi ý văn bản thành tham số của mô hình, cho phép quá trình suy diễn loại bỏ hoàn toàn gánh nặng tính toán của bộ mã hóa ngôn ngữ. Điều này đưa tốc độ xử lý về mức tương đương mạng CNN truyền thống.

Hình~\ref{fig:vis_dense} minh họa một trường hợp đếm thực tế với mật độ dày đặc (175 đối tượng): negative exemplar (đỏ, góc trên trái) đánh dấu vùng nền, positive exemplars (xanh lá) chọn đúng marker pen. Dù mô hình đếm thiếu đáng kể (dự đoán 124,89), ví dụ này cho thấy rõ vai trò của NSM trong việc phân biệt foreground/background.

\begin{figure}[H]
    \centering
    \figureorplaceholder{figures/vis_dense_negexemplar.jpg}{0.95\textwidth}{Minh họa VA-Count trên ảnh dày đặc \texttt{6822.jpg}: positive exemplars (xanh lá) và negative exemplar (đỏ). GT=175, Predicted=124.89. Trường hợp đếm thiếu do mật độ quá dày.}
    \caption{Trường hợp mật độ dày đặc (\texttt{6822.jpg}, lớp \emph{marker pens}): positive exemplars (xanh lá) và negative exemplar (đỏ, góc trên trái). GT $= 175$, dự đoán $\approx 124{,}89$. Đây là ví dụ điển hình về lỗi đếm thiếu do vật thể quá nhỏ và sát nhau.}
    \label{fig:vis_dense}
\end{figure}

\begin{figure}[H]
    \centering
    \figureorplaceholder{figures/yolo_speed_curve.png}{0.65\textwidth}{Speed-and-Accuracy Curve so sánh YOLO-World với GroundingDINO/GLIP trên tập LVIS minival (trích từ~\cite{yoloworld}).}
    \caption{Speed-and-Accuracy Curve: YOLO-World so với GroundingDINO/GLIP trên LVIS minival. Đồ thị trích từ bài báo gốc YOLO-World~\cite{yoloworld}.}
    \label{fig:yolo_speed}
\end{figure}

YOLO-World đạt 35 mAP với 50--80 FPS trên LVIS minival, so với 27 mAP và <5 FPS của GroundingDINO-T. Khi tích hợp YOLO-World vào VA-Count, hai luồng EEM được thay GroundingDINO tương ứng (\texttt{yolo\_pos\_withPrompt.py}, \texttt{yolo\_pos\_withoutPrompt.py}, \texttt{yolo\_neg.py}), trong khi NSM giữ nguyên.

\begin{figure}[H]
    \centering
    \figureorplaceholder{figures/system_yolo.png}{0.95\textwidth}{Kiến trúc VA-Count với YOLO-World thay thế GroundingDINO. Luồng prompt enhancement vẫn được dùng tùy chọn.}
    \caption{Kiến trúc hệ thống sau khi thay GroundingDINO bằng YOLO-World ở module phát hiện. Sơ đồ chuyển vẽ theo thiết kế của VA-Count~\cite{vacount} với phần phát hiện được thay theo YOLO-World~\cite{yoloworld}.}
    \label{fig:system_yolo}
\end{figure}

%==================================================================
%                          CHƯƠNG 3
%==================================================================
\chapter{THỰC NGHIỆM VÀ KẾT QUẢ}
\label{ch:experiments}

\section{Bộ dữ liệu FSC147}

FSC147~\cite{ranjan2021fsc147} là dataset chuẩn cho \textit{class-agnostic / zero-shot counting}: 6\,135 ảnh, 147 lớp, mỗi ảnh được chú thích các \textit{dot annotations}. Bộ dữ liệu nổi bật với:
\begin{itemize}
    \item Các tập con không chồng lấn (train/val/test theo lớp).
    \item Độ thách thức cao do biến thiên lớn về mật độ và kích thước đối tượng.
    \item Nhãn lớp + chú thích dạng điểm cho phép đánh giá zero-shot thông qua text prompts.
\end{itemize}

Nhóm áp dụng chiến lược chia 6:2:2 (train:val:test) theo split chính thức của FSC-147~\cite{ranjan2021fsc147}.

\begin{figure}[H]
    \centering
    \figureorplaceholder{figures/fsc147_samples.png}{0.85\textwidth}{Một số ảnh minh họa từ FSC147: từ các lớp dày đặc (oranges, beans, keyboard keys) đến các lớp thưa.}
    \caption{Một vài ảnh minh họa từ bộ dữ liệu FSC-147~\cite{ranjan2021fsc147}.}
    \label{fig:fsc147}
\end{figure}

Hình~\ref{fig:vis_good} và~\ref{fig:vis_peas} minh họa kết quả đếm định tính trên hai ảnh từ tập test. Hệ thống cho thấy khả năng đếm chính xác trên nhiều loại đối tượng đa dạng, với density map tập trung đúng vào vị trí các đối tượng.

\begin{figure}[H]
    \centering
    \figureorplaceholder{figures/vis_good_prediction.png}{0.95\textwidth}{Kết quả tốt: ảnh 4910.jpg (chocolates), GT=14, Predicted=14.40, sai số 2.86\%.}
    \caption{Kết quả đếm tốt trên ảnh \texttt{4910.jpg} (lớp \emph{chocolates}): 3 positive exemplars được chọn đúng, density map phân bố chính xác. GT $= 14$, dự đoán $= 14{,}40$ (sai số $2{,}86\%$).}
    \label{fig:vis_good}
\end{figure}

\begin{figure}[H]
    \centering
    \figureorplaceholder{figures/vis_peas_result.png}{0.95\textwidth}{Kết quả đếm trên ảnh 5511.jpg (peas), GT=18, Predicted=19.48, sai số 8.22\%.}
    \caption{Kết quả đếm trên ảnh \texttt{5511.jpg} (lớp \emph{peas}): dù đối tượng nhỏ và xếp dày, hệ thống vẫn duy trì sai số thấp. GT $= 18$, dự đoán $= 19{,}48$ (sai số $8{,}22\%$).}
    \label{fig:vis_peas}
\end{figure}

\section{Độ đo đánh giá}

Hệ thống được đánh giá bằng hai độ đo chuẩn cho counting: MAE và RMSE.

\subsection{Mean Absolute Error (MAE)}
MAE đo trung bình giá trị sai số tuyệt đối:
\begin{equation}
\mathrm{MAE} = \frac{1}{N}\sum_{i=1}^{N}\left|y_i - \hat{y}_i\right|,
\end{equation}
trong đó $N$ là số ảnh, $y_i$ là số đối tượng thực tế (ground truth), $\hat{y}_i$ là số đếm dự đoán.

\subsection{Root Mean Squared Error (RMSE)}
RMSE phạt nặng các sai số lớn nhờ bình phương trước khi lấy căn:
\begin{equation}
\mathrm{RMSE} = \sqrt{\frac{1}{N}\sum_{i=1}^{N}(y_i - \hat{y}_i)^2}.
\end{equation}

Sử dụng đồng thời cả hai cho phép đánh giá toàn diện: MAE phản ánh độ chính xác trung bình; RMSE phản ánh độ ổn định trên các trường hợp khó.

\section{Thiết lập thực nghiệm}
\label{sec:setup}

Các con số trong các bảng dưới đây được thu thập trên cấu hình:
\begin{itemize}
    \item \textbf{Huấn luyện và đánh giá:} server NVIDIA H200, 24 nhân CPU, 192\,GB RAM, 3\,TB NVMe.
    \item \textbf{Demo:} laptop Intel Core i7-13620H + RTX~4060.
\end{itemize}

\begin{table}[H]
\centering
\renewcommand{\arraystretch}{1.2}
\begin{tabular}{|l|c|}
\hline
\rowcolor[HTML]{EFEFEF}
\textbf{Tham số} & \textbf{Giá trị} \\
\hline
Optimizer & AdamW \\
\hline
Base Learning Rate & $1{\times}10^{-5}$ \\
\hline
Weight Decay & $0.05$ \\
\hline
LR Schedule & Cosine Annealing \\
\hline
Total Epochs & $500$ \\
\hline
Batch Size (per GPU) & $8$ \\
\hline
Gradient Accumulation & $16$ steps \\
\hline
\end{tabular}
\caption{Siêu tham số huấn luyện được nhóm sử dụng, theo cấu hình khuyến nghị trong VA-Count~\cite{vacount}.}
\label{tab:hyperparams}
\end{table}

\paragraph{Lưu ý.}
CLI mặc định trong \texttt{FSC\_train.py} có giá trị khác
(epochs $1000$, lr $5{\times}10^{-5}$).
Bảng~\ref{tab:hyperparams} liệt kê tổ hợp đã được nhóm sử dụng cho các
kết quả báo cáo, lưu trong log \texttt{wandb} ở thư mục
\texttt{experiments/exp\{2,3,4,5\}/wandb/}.

\section{Kết quả thực nghiệm}
\label{sec:results}

Các kết quả ở Bảng~\ref{tab:extract_time}, Bảng~\ref{tab:counting_perf},
Bảng~\ref{tab:demo_time} được tổng hợp từ các log \texttt{wandb} của nhóm
trong \mbox{\texttt{experiments/exp\{2,3,4,5\}/wandb/}}
và đã được đối chiếu lại với các tệp \texttt{wandb-summary.json} còn lưu được.

\subsection{Thời gian trích xuất exemplars}

\begin{table}[H]
\centering
\renewcommand{\arraystretch}{1.2}
\begin{tabular}{|l|c|c|}
\hline
\rowcolor[HTML]{EFEFEF}
\textbf{Phương pháp} & \textbf{Positive (giờ)} & \textbf{Negative (giờ)} \\
\hline
Grounding DINO (Baseline) & $\approx 3$ & $\approx 10$ \\
\hline
Grounding DINO + Rich Prompt & $\approx 10$ & $\approx 15$ \\
\hline
YOLO-World & $\approx 0.5$ & $\approx 2$ \\
\hline
YOLO-World + Rich Prompt & $\approx 1$ & $\approx 3$ \\
\hline
\end{tabular}
\caption{Thời gian trích xuất exemplars trên toàn bộ 6\,146 ảnh FSC147 (đo bởi nhóm trên một workstation RTX~4060 đơn).}
\label{tab:extract_time}
\end{table}

YOLO-World hoàn thành trích xuất positive boxes chỉ trong $\approx 0{.}5$ giờ, nhanh gấp 6 lần GroundingDINO ($\approx 3$ giờ). Khi tích hợp Rich Prompt, ưu thế càng rõ: GroundingDINO + Rich Prompt mất tới 25 giờ tổng, trong khi YOLO-World + Rich Prompt chỉ tốn 4 giờ. Đây là một khác biệt mang tính quyết định cho việc chuẩn bị dữ liệu trên dataset lớn.

\subsection{Hiệu suất đếm trên FSC-147}

\begin{table}[H]
\centering
\renewcommand{\arraystretch}{1.2}
\begin{tabular}{|l|c|c|}
\hline
\rowcolor[HTML]{EFEFEF}
\textbf{Mô hình} & \textbf{MAE} $\downarrow$ & \textbf{RMSE} $\downarrow$ \\
\hline
VA-Count & 17,99 & 129,39 \\
\hline
VA-Count + Rich Prompt & \textbf{17,80} & 129,69 \\
\hline
VA-Count với YOLO-World & 19,03 & 131,55 \\
\hline
VA-Count với YOLO-World + Rich Prompt & 17,91 & 130,98 \\
\hline
\end{tabular}
\caption{So sánh hiệu suất đếm trên FSC-147 (test split). Các giá trị MAE/RMSE được tổng hợp từ log \texttt{wandb} của nhóm trong thư mục \mbox{\texttt{experiments/exp\{2,3,4,5\}/wandb/}}.}
\label{tab:counting_perf}
\end{table}

\paragraph{Phân tích.}
\begin{itemize}
    \item \textbf{Rich Prompt nâng nhẹ chất lượng đếm:} MAE giảm từ 17,99 xuống 17,80 trên VA-Count. Trên YOLO-World, mức cải thiện rõ hơn (19,03 $\to$ 17,91).
    \item \textbf{YOLO-World một mình hơi xấu hơn baseline} (MAE 19,03 vs 17,99) do nhiễu phát hiện cao hơn ở giai đoạn đếm.
    \item \textbf{YOLO-World + Rich Prompt là điểm cân bằng tốt nhất:} gần như ngang VA-Count gốc về độ chính xác (MAE 17,91), nhưng chi phí tính toán thấp hơn rất nhiều (xem Bảng~\ref{tab:demo_time}).
\end{itemize}

\subsection{Tốc độ inference khi demo}

\begin{table}[H]
\centering
\renewcommand{\arraystretch}{1.2}
\begin{tabular}{|l|c|}
\hline
\rowcolor[HTML]{EFEFEF}
\textbf{Phương pháp} & \textbf{Thời gian trung bình (s/ảnh)} \\
\hline
VA-Count & 1,4710 \\
\hline
VA-Count + Rich Prompt & 5,7578 \\
\hline
VA-Count với YOLO-World & \textbf{0,6006} \\
\hline
VA-Count với YOLO-World + Rich Prompt & 2,4054 \\
\hline
\end{tabular}
\caption{Thời gian inference trung bình trên 1 ảnh trong môi trường demo (RTX~4060).}
\label{tab:demo_time}
\end{table}

Trên cấu hình demo, YOLO-World giảm thời gian xuống $0{,}60$ s/ảnh, nhanh gấp 2,5 lần VA-Count gốc. Khi bật Rich Prompt, hệ thống giữ ở mức $2{,}4$ s/ảnh, vẫn chấp nhận được với tương tác người dùng. Phương án \textbf{YOLO-World + Rich Prompt} đạt cân bằng tốt nhất giữa độ chính xác (MAE $17{,}91$) và độ trễ ($2{,}4$ s/ảnh), và đây là phương án mặc định mà ứng dụng demo của nhóm sử dụng (Chương~\ref{ch:demo}).

\section{Phân tích lỗi sai}
\label{sec:error-analysis}

CS338.Q21 phân loại các lỗi quan sát được trên FSC-147 thành ba nhóm:

\subsection{Mật độ dày đặc}

Khi đối tượng nhỏ và sát nhau (ví dụ hộp bút màu hoặc hạt cườm), ranh giới thị giác mờ nhạt, làm density map suy giảm giá trị và gây \textbf{đếm thiếu}.

\begin{figure}[H]
    \centering
    \figureorplaceholder{figures/failure_dense.png}{0.95\textwidth}{Trường hợp đếm thiếu khi mật độ dày đặc.}
    \caption{Minh họa sai số đếm thiếu trên ảnh có mật độ rất dày (ảnh thực tế từ \texttt{experiments/exp3/visualizations\_test/worst\_1\_687.png}, GT $= 548$, dự đoán $\approx 283{,}55$, sai số tuyệt đối $264{,}45$).}
    \label{fig:fail_dense}
\end{figure}

\subsection{Phân mảnh đối tượng}

Các vật thể có cấu trúc chia tách (cửa sổ nhiều ô, kính mắt) bị nhận diện nhầm thành nhiều thực thể độc lập, tạo nhiều cực đại cục bộ trên một đối tượng và gây \textbf{đếm thừa}.

\begin{figure}[H]
    \centering
    \figureorplaceholder{figures/failure_fragment.png}{0.85\textwidth}{Trường hợp đếm thừa do phân mảnh đối tượng (cửa sổ nhiều ô).}
    \caption{Failure case: phân mảnh đối tượng. Dãy cửa sổ nhiều ô bị tách thành nhiều cực đại cục bộ trên density map (GT $= 13$, dự đoán $\approx 27{,}13$). Ví dụ định tính trích từ phân tích nội bộ của nhóm trên FSC-147.}
    \label{fig:fail_fragment}
\end{figure}

\subsection{Sai số mang tính phương pháp}

Mô hình dùng một bộ trích xuất kiểu Masked Autoencoder + density map. Quan sát thực tế cho thấy mô hình \emph{vẫn đưa ra ước lượng} ngay cả khi không lấy được exemplar hợp lệ, tuy nhiên sai số rất lớn (ví dụ ground truth $371$, dự đoán $251{,}55$). Điều này gợi ý mô hình có xu hướng dựa vào đặc trưng toàn cục (mật độ, kết cấu, bố cục không gian) thay vì điều kiện hóa thực sự theo exemplar. Đây là một \textbf{bias dữ liệu / kiến trúc} cần được giải quyết trong các phiên bản kế tiếp (xem Chương~\ref{ch:conclusion}).

%==================================================================
%                          CHƯƠNG 4
%==================================================================
\chapter{DEMO VÀ ĐÁNH GIÁ HỆ THỐNG}
\label{ch:demo}

\begin{flushleft}
    \textbf{Mã nguồn dự án:} \href{https://github.com/paht2005/CS338.Q21_Zero-shot-Object-Coutning-with-Good-Examplers}{GitHub Repository} \\
    \textbf{Video Demo:} cập nhật trong file \texttt{README.md} (mục \emph{Demo \& Outputs}).
\end{flushleft}

Nhóm cung cấp một ứng dụng demo dạng Streamlit~\cite{streamlit} (\texttt{code/source-code/demo\_app\_advanced.py}) cùng các script hỗ trợ:

\begin{itemize}
    \item \texttt{demo\_app\_advanced.py}: ứng dụng Streamlit nâng cao, cho phép tùy chọn extractor (GroundingDINO / YOLO-World), bật/tắt Rich Prompt, hiển thị bản đồ mật độ, bounding box positive/negative và so sánh với ground-truth (mặc định: \textbf{YOLO-World + Rich Prompt}).
    \item \texttt{demo\_inference.py}: thư viện inference dùng chung (load model, expand prompts, pipeline tối giản); được \texttt{demo\_app\_advanced.py} import.
    \item \texttt{demo\_pipeline\_advanced.py}, \texttt{demo\_visualization.py}: module pipeline đầy đủ và các hàm vẽ overlay được dùng bên trong ứng dụng Streamlit.
\end{itemize}

\section{Cài đặt và chạy demo}

\begin{lstlisting}[style=bashstyle]
cd code/source-code
python -m venv .venv && source .venv/bin/activate
pip install -r requirements.txt

# 1. Set Gemini API key (DO NOT commit)
cp .env.example .env
# edit .env and fill in GEMINI_API_KEY

# 2. Run the Streamlit demo
streamlit run demo_app_advanced.py
\end{lstlisting}

\section{Giao diện và tương tác}
Người dùng tải lên một ảnh tự nhiên, nhập tên lớp tiếng Anh (ví dụ \texttt{oranges}), và nhận về:
\begin{itemize}
    \item Tổng số đối tượng đếm được.
    \item Bản đồ mật độ overlay trên ảnh gốc.
    \item Các bounding box positive/negative đã được EEM chọn.
\end{itemize}

\begin{figure}[H]
    \centering
    \figureorplaceholder{figures/demo_streamlit.png}{0.95\textwidth}{Đầu ra ba ô (ảnh gốc với positive exemplars, predicted density map, overlay) cùng bảng so sánh số lượng dự đoán và ground truth.}
    \caption{Đầu ra của hệ thống cho ảnh \texttt{346.jpg} (lớp \emph{strawberries}): ảnh gốc kèm 3 positive exemplars (trái), predicted density map (giữa), overlay heatmap lên ảnh gốc (phải) và bảng tổng hợp số đếm. Ground truth $= 11$, mô hình dự đoán $\approx 11{,}26$ (sai số tuyệt đối $0{,}26$). Visualization từ \texttt{experiments/exp3/visualizations\_test/best\_1\_346.png}; ứng dụng demo Streamlit dùng cùng pipeline để render.}
    \label{fig:demo}
\end{figure}

\section{Đánh giá định tính}
Kiểm thử trên các ảnh khó từ FSC-147 và một số ảnh ngoài tập (chợ trái cây, ảnh chăn nuôi) cho thấy:
\begin{itemize}
    \item Cấu hình \textbf{YOLO-World + Rich Prompt} cho phản hồi nhanh ($\sim 2$ s) và chấp nhận được trên CPU/GPU laptop.
    \item Khi class name rõ ràng và đối tượng phân tách tốt, kết quả gần với ground truth.
    \item Trên các ảnh dày đặc hoặc class name mơ hồ, hệ thống bộc lộ đúng các lỗi đã phân loại ở Mục~\ref{sec:error-analysis}.
\end{itemize}

%==================================================================
%                          CHƯƠNG 5
%==================================================================
\chapter{KẾT LUẬN VÀ HƯỚNG PHÁT TRIỂN}
\label{ch:conclusion}

\section{Kết luận}
Đồ án này hiện thực và đánh giá một hệ thống Zero-shot Object Counting dựa trên VA-Count~\cite{vacount}, mở rộng theo Rich Prompts~\cite{richprompts} và YOLO-World~\cite{yoloworld}. Các kết luận chính:

\begin{itemize}
    \item \textbf{Về kiến trúc.} VA-Count là một thiết kế giàu thông tin học được tách rõ giữa pha lấy mẫu (EEM) và pha sinh density map (NSM). Việc chia luồng dương/âm tạo điều kiện cho cơ chế contrastive, một nguyên lý vẫn hữu hiệu cho zero-shot counting.
    \item \textbf{Về cải tiến.} Rich Prompt giảm MAE đáng kể khi kết hợp với YOLO-World ($19{,}03 \to 17{,}91$) và nhẹ trên VA-Count baseline ($17{,}99 \to 17{,}80$), khẳng định vai trò của ngữ nghĩa chi tiết cho mô hình mở-từ-vựng. YOLO-World loại bỏ được nút thắt cổ chai về tốc độ (giảm thời gian extraction từ $\sim 25h$ xuống $\sim 4h$) và đem lại độ trễ inference chấp nhận được cho demo.
    \item \textbf{Về cài đặt.} Nhóm dùng dạng \emph{margin-ReLU} cho hàm contrastive thay vì dạng softmax như Công thức (2) trong~\cite{vacount} (tương đương về tối ưu, ổn định số học hơn); dùng CLIP ViT-L/14 cho re-ranking patch--text bên cạnh CLIP ViT-B/32 cho bộ phân loại nhị phân; và dùng Gemini~2.0~Flash (\texttt{gemini-2.0-flash-exp}) thay cho phiên bản 2.5 Flash trong bài báo Rich Prompts.
    \item \textbf{Về kỹ thuật phần mềm.} Khóa Gemini API được lưu trong \texttt{.env} (không commit), kho mã nguồn được tổ chức theo \texttt{code/}, \texttt{experiments/}, \texttt{docs/}, \texttt{configs/}, \texttt{scripts/} và đi kèm với hai pin file (\texttt{requirements.txt} mặc định, \texttt{requirements-cuda116.txt} cho tái lập số liệu).
\end{itemize}

\section{Hướng phát triển}

Trên cơ sở các phân tích lỗi sai (Mục~\ref{sec:error-analysis}) và các lựa chọn cài đặt nêu trên, nhóm đề xuất các hướng phát triển sau:

\begin{itemize}
    \item \textbf{Khử phân mảnh đối tượng.} Áp dụng các kỹ thuật post-processing (graph-based clustering, watershed) hoặc thêm head phụ dự đoán \textit{instance-aware density} để giảm đếm thừa cho các vật thể có cấu trúc chia tách.
    \item \textbf{Mật độ dày đặc.} Bổ sung kiến trúc đa độ phân giải hoặc tile-based inference để giữ chi tiết tại biên các vật thể sát nhau.
    \item \textbf{Khử bias đếm-không-cần-exemplar.} Thử nghiệm các phương pháp đếm dựa trên matching trực tiếp (CountR~\cite{countr}) thay cho density map MAE-based, hoặc thêm regularizer ép mô hình bằng 0 khi exemplar không hợp lệ.
    \item \textbf{Khảo sát các lựa chọn cài đặt.} Thử nghiệm dạng softmax contrastive theo Công thức (2) trong~\cite{vacount} và so sánh với dạng margin-ReLU; so sánh CLIP ViT-B/32 với ViT-L/14 cho re-ranking; chạy lại Rich Prompt với Gemini~2.5~Flash khi có quota.
    \item \textbf{Triển khai sản phẩm.} Đóng gói pipeline trong Docker và triển khai trên dịch vụ HuggingFace Spaces / Vercel để tăng khả năng tiếp cận.
\end{itemize}

%==================================================================
%                          BIBLIOGRAPHY
%==================================================================

\addcontentsline{toc}{chapter}{TÀI LIỆU THAM KHẢO}

\begin{thebibliography}{99}

\bibitem{vacount}
H.~Zhu, J.~Yuan, Z.~Yang, Y.~Guo, Z.~Wang, X.~Zhong, and S.~He,
``Zero-shot object counting with good exemplars,''
in \emph{European Conference on Computer Vision (ECCV)}, 2024.
[Online]. Available:
\url{https://www.ecva.net/papers/eccv_2024/papers_ECCV/papers/00812.pdf}

\bibitem{richprompts}
H.~Zhu, S.~Li, J.~Yuan, Z.~Yang, Y.~Guo, W.~Liu, X.~Zhong, and S.~He,
``Expanding zero-shot object counting with rich prompts,''
arXiv preprint arXiv:2505.15398, 2025.
[Online]. Available:
\url{https://arxiv.org/abs/2505.15398}

\bibitem{countr}
C.~Liu, Y.~Zhong, A.~Zisserman, and W.~Xie,
``CounTR: Transformer-based generalised visual counting,''
in \emph{British Machine Vision Conference (BMVC)}, 2022.
[Online]. Available:
\url{https://pdf-files.bmvc2022.org/0370.pdf}

\bibitem{yoloworld}
T.~Cheng, L.~Song, Y.~Ge, W.~Liu, X.~Wang, and Y.~Shan,
``YOLO-World: Real-time open-vocabulary object detection,''
in \emph{IEEE/CVF Conference on Computer Vision and Pattern
Recognition (CVPR)}, 2024.
[Online]. Available:
\url{https://openaccess.thecvf.com/content/CVPR2024/papers/Cheng_YOLO-World_Real-Time_Open-Vocabulary_Object_Detection_CVPR_2024_paper.pdf}

\bibitem{groundingdino}
S.~Liu \emph{et~al.},
``Grounding DINO: Marrying DINO with grounded pre-training
for open-set object detection,''
in \emph{European Conference on Computer Vision (ECCV)}, 2024.
[Online]. Available:
\url{https://arxiv.org/abs/2303.05499}

\bibitem{clip}
A.~Radford \emph{et~al.},
``Learning transferable visual models from natural language
supervision (CLIP),''
in \emph{International Conference on Machine Learning (ICML)},
2021.
[Online]. Available:
\url{https://arxiv.org/abs/2103.00020}

\bibitem{ranjan2021fsc147}
V.~Ranjan, U.~Sharma, T.~Nguyen, and M.~Hoai,
``Learning to count everything (FSC-147),''
in \emph{IEEE/CVF Conference on Computer Vision and Pattern
Recognition (CVPR)}, 2021.
[Online]. Available:
\url{https://arxiv.org/abs/2104.08391}

\bibitem{gemini}
Google DeepMind, ``Gemini API documentation,'' 2024--2026.
[Online]. Available:
\url{https://ai.google.dev/gemini-api/docs}

\bibitem{streamlit}
Streamlit, ``The fastest way to build and share data apps,'' 2026.
[Online]. Available: \url{https://streamlit.io}

\bibitem{mae}
K.~He, X.~Chen, S.~Xie, Y.~Li, P.~Doll{\'a}r, and R.~Girshick,
``Masked autoencoders are scalable vision learners,''
in \emph{IEEE/CVF Conference on Computer Vision and Pattern
Recognition (CVPR)}, 2022.
[Online]. Available:
\url{https://arxiv.org/abs/2111.06377}

\bibitem{famnet}
M.~Ranasinghe, S.~Herath, H.~Bilen, and F.~Zheng,
``Learning to count without annotations,''
in \emph{IEEE/CVF Conference on Computer Vision and Pattern
Recognition (CVPR)}, 2022.
[Online]. Available:
\url{https://arxiv.org/abs/2307.08727}

\bibitem{bmnet}
M.~Shi, H.~Lu, C.~Feng, C.~Liu, and Z.~Cao,
``Represent, compare, and learn: A similarity-aware framework
for class-agnostic counting,''
in \emph{IEEE/CVF Conference on Computer Vision and Pattern
Recognition (CVPR)}, 2022.
[Online]. Available:
\url{https://arxiv.org/abs/2203.08354}

\bibitem{simclr}
T.~Chen, S.~Kornblith, M.~Norouzi, and G.~Hinton,
``A simple framework for contrastive learning of visual
representations (SimCLR),''
in \emph{International Conference on Machine Learning (ICML)}, 2020.
[Online]. Available:
\url{https://arxiv.org/abs/2002.05709}

\bibitem{zeroshot_counting_xu}
V.~Xu, L.~Zhu, and Y.~Yang,
``Zero-shot object counting,''
in \emph{IEEE/CVF Conference on Computer Vision and Pattern
Recognition (CVPR)}, 2023.
[Online]. Available:
\url{https://arxiv.org/abs/2303.02001}

\bibitem{ovdet_survey}
Y.~Zang, W.~Li, J.~Han, K.~Zhou, and C.~C.~Loy,
``Open-vocabulary object detection: A survey,''
arXiv preprint arXiv:2312.15850, 2023.
[Online]. Available:
\url{https://arxiv.org/abs/2312.15850}

\end{thebibliography}

\end{document}
